% Artículo VII. Incorporación de resultados preexistentes, 10-09-2026.
% Propietarios íntegros preservados en fuentes_conservadas:
% 17_pantalla_cubica_soldadura.tex (P), 18_clausura_energetica.tex (E).
% P: incidencia, cociente, soldadura celular y refinamiento.
% E: respuesta, unicidad, extensión mínima y expectativa tracial.
% Dependencias anteriores de VII: núcleo APP--TRIT--TPK; realización
% cúbica de la ruta, cociclo de transporte y compresión q_vis=5/9.
% Este fragmento no constituye por sí solo el artículo autónomo.
% La identidad de pantalla se distingue de la contracción material de espín.

\subsection{Incidencia cúbica y soldadura de frontera}
\label{vii:sec:pantalla-respuesta}

La realización cúbica del transporte conserva seis incidencias orientadas.
Su cociente geométrico, la forma lorentziana y la respuesta energética
se obtienen de la misma matriz de incidencias. El transporte procede de
la ruta enriquecida APP--TRIT--TPK: la cara, el signo, el acarreo y la
memoria acompañan a cada arista. En esta sección se desarrolla la
operación lineal que actúa sobre esa representación de frontera.

\subsubsection{Incidencia cúbica y cociente de rango cuatro}
\label{vii:sec:cociente-pantalla}

Sea \(H_F=\mathbb R^F\), con
\(F=\{(i,\epsilon):1\leq i\leq3,\ \epsilon\in\{-1,1\}\}\),
su producto escalar euclidiano y su base \(e_{i,\epsilon}\).
La oposición de caras es la involución ortogonal
\(Re_{i,\epsilon}=e_{i,-\epsilon}\). Definimos
\[
 u=\sum_{i,\epsilon}e_{i,\epsilon},\qquad
 P_u=\frac{uu^{\mathsf T}}6,\qquad P_-=\frac{I-R}{2},\qquad
 P_0=\frac{I+R}{2}-P_u,\qquad P_\square=P_u+P_-.
\]
Los vértices son \(V=\{-1,1\}^3\). La matriz de incidencia
\(M:\mathbb R^V\longrightarrow H_F\) queda determinada por
\[
 M_{(i,\epsilon),\nu}=\mathbf1_{\{\nu_i=\epsilon\}}.
\]

\begin{proposition}[Descomposición de la incidencia]
\label{vii:prop:incidencia}
Se tiene
\[
 MM^{\mathsf T}=12P_u+4P_-,
 \qquad
 \ker M^{\mathsf T}=P_0H_F.
\]
Por consiguiente, el cociente
\(Q_\square=Q=H_F/P_0H_F\), identificado ortogonalmente con
\(P_\square H_F\), tiene dimensión cuatro.
\end{proposition}

\begin{proof}
Una cara contiene cuatro vértices. Dos caras opuestas carecen de vértices
comunes y dos caras de ejes diferentes tienen dos vértices comunes.
Escribiendo \(\mathbf J=uu^{\mathsf T}\), resulta
\[
 MM^{\mathsf T}=4I+2(\mathbf J-I-R)
 =2(I-R)+2\mathbf J=4P_-+12P_u.
\]
Los tres proyectores \(P_u,P_-,P_0\) son ortogonales y suman \(I\);
sus rangos son, respectivamente, uno, tres y dos. La identidad
\(\|M^{\mathsf T}x\|^2=\langle x,MM^{\mathsf T}x\rangle\)
identifica el núcleo anunciado. El rango complementario es cuatro.
\end{proof}

Sobre \(Q\), la forma
\[
 B=P_--P_u
\]
tiene signatura \((3,1)\). La orientación temporal se fija por \(u\).
Con
\[
 t=\frac{u}{\sqrt6},\qquad
 d_i=\frac{e_{i,+}-e_{i,-}}{\sqrt2},\qquad
 W=\sqrt6P_u+\sqrt2P_-,
\]
se cumple \(We_{i,\epsilon}=t+\epsilon d_i\). Esos seis vectores
son nulos para \(B\), mientras que
\[
 3t+\nu_1d_1+\nu_2d_2+\nu_3d_3
\]
tiene norma cuadrática \(-6\). La misma incidencia satisface
\(MM^{\mathsf T}=2W^*W\) sobre \(Q\). Las rotaciones propias del cubo
conmutan con los proyectores, preservan ambas formas y fijan \(t\).
Estas identidades determinan la métrica de la representación cúbica.

\subsubsection{Soldadura celular y transporte de memoria}
\label{vii:sec:soldadura-celular}

Para cada arista orientada \(a\), la lectura de la ruta proporciona
la cara \(\lambda(\underline a)\), el signo
\(\sigma(a)\), las fibras de origen y destino y el transporte
\(U(a):Q_{s(a)}\longrightarrow Q_{t(a)}\).
La inversión verifica \(U(\bar a)=U(a)^{-1}\).
Las semiaristas de frontera conservan su inversión formal y su incidencia
en el borde. Con \(\ell_m=\ell_0\,9^{-m}>0\), la soldadura es
\begin{equation}
 \beta_m(a)=\ell_m\sigma_m(a)
              n_{\lambda(\underline a)}^{s(a)},\qquad
 n_{i,\epsilon}=t+\epsilon d_i.
 \label{vii:eq:soldadura-arista}
\end{equation}
La cara y el signo se transportan junto con la memoria de la ruta.
La convención de inversión exige
\begin{equation}
 \beta_m(\bar a)=-U_m(a)\beta_m(a).
 \label{vii:eq:inversion-soldadura}
\end{equation}
Aplicando dos veces esta relación se recupera \(\beta_m(a)\);
el registro de acarreo conserva, por separado, la composición ordenada.

\begin{proposition}[Naturalidad de la soldadura por refinamiento]
\label{vii:prop:refinamiento-soldadura}
Sea \(a=a_1\cdots a_9\) un refinamiento compatible de una arista.
Supóngase que las nueve incidencias hijas, transportadas a la fibra del
predecesor, conservan la cara y el signo, y que
\(\ell_{m+1}=\ell_m/9\). Sea
\(\mathcal U_j:Q_{m,s(a)}\longrightarrow Q_{m+1,s(a_j)}\)
el transporte acumulado que incluye la identificación por refinamiento
y el trayecto hasta el origen del sucesor \(j\). Entonces
\[
 \beta_m(a)=
 \sum_{j=1}^{9}\mathcal U_j^{-1}\beta_{m+1}(a_j).
\]
La igualdad es compatible con la inversión y con la concatenación
del registro de memoria.
\end{proposition}

\begin{proof}
Cada sumando, expresado en la fibra del origen del predecesor, es
\((\ell_m/9)\sigma_m(a)n_{\lambda(\underline a)}\).
La suma de nueve términos reproduce la soldadura del predecesor.
La inversión invierte el orden de los transportes y sustituye cada uno
por su inverso; la relación \eqref{vii:eq:inversion-soldadura} proporciona
el signo global. La memoria se concatena en ese mismo orden, por lo
que el refinamiento conserva tanto la representación vectorial como su
registro enriquecido.
\end{proof}
